% 2de graad, Reële getallen (getallen, machten en vierkantswortels).
% Van Doesburg: een raster dat over zijn diagonaal gekanteld is.
% Motief: de diagonaal van het eenheidsvierkant, met de passer op de getallenas gelegd: daar ligt sqrt(2).
\def\omslagbestand{omslag-contracompositie.jpg}
\def\omslagbron{Theo van Doesburg, \textit{Contra-compositie VI}, 1925 (publiek domein)}
\def\omslagrand{10.6}
\def\omslagmotief{%
  \randlijn
  \draw[omslaglijn=1pt] (14,10.6) rectangle (16,12.6);
  \draw[omslaglijn=.8pt,dashed] (16,12.6) arc[start angle=45,end angle=0,radius=2.8284];
  \draw[fel,line width=2pt] (14,10.6) -- (16,12.6);
  \foreach \i in {0,1,2,3}{%
    \draw[white,line width=1.2pt] ({14+2*\i},10.35) -- ({14+2*\i},10.85);
    \node[text=white,font=\small,anchor=north] at ({14+2*\i},10.25) {$\i$};
  }
  \fill[omslagstip] (16.8284,10.6) circle (1.4mm);
  \node[omslaglabel,anchor=south west] at (16.95,10.75) {$\sqrt2$};
}
